\documentclass[10pt,a4paper]{amsart}
\usepackage[margin=19mm]{geometry}
\usepackage{amsmath,amssymb,mathtools}
\usepackage[scr=rsfs,cal=euler]{mathalfa}
\usepackage{xfrac}
\usepackage[unicode,hidelinks,bookmarks=false]{hyperref}
\newcommand{\ie}{i.e.\ }
\newcommand{\cf}{cf.\ }
\newcommand{\resp}{respectively\ }
\newcommand{\Subsection}{\subsection}
\providecommand{\nobreakdash}{\nobreak\hbox{-}}
\setlength{\emergencystretch}{2em}
\multlinegap0pt
\usepackage{tikz}
\usepackage{amssymb}
\usepackage[shortlabels]{enumitem}
\newtheorem{proposition}{Proposition}
\title{Geometric Invariant Theory of Peterson Varieties}
\author{Arkadev Ghosh, Santosha Pattanayak}
\date{Source: arXiv:2608.26695v1 (CC BY 4.0)}
\begin{document}
\maketitle

\noindent\emph{Source and licence.} Geometric Invariant Theory of Peterson Varieties. Version arXiv:2608.26695v1 (CC BY 4.0), distributed by arXiv under the Creative Commons Attribution 4.0 International licence, http://creativecommons.org/licenses/by/4.0/, as recorded in the arXiv licence metadata for this exact version and shown on the abstract page https://arxiv.org/abs/2608.26695v1. Full licence notice: \texttt{licenses/arxiv-2608.26695-CC-BY-4.0.txt}.\par\smallskip

\noindent\textit{Editorial scope.} Counted unit: the article's proposition prop:maximal-stratum-orbits (source0.tex lines 2390--2416). The article's complete original proof of it is delivered with it (source0.tex lines 2418--2618, one \texttt{proof} environment of 201 lines). No mathematics is written, repaired or re-derived: the statement and the proof are the article's own lines, in the article's own notation, and no retained line is substituted or re-flowed.  No further source-local context is needed: every cross-reference of every retained line resolves inside the two delivered spans.  Nothing was omitted from any retained span, so the package carries no omission record.  No retained line contains a citation, so the wrapper carries no bibliography and no bibliography entry.  No retained line contains a figure environment, an image-inclusion command or drawing code, so no image is delivered and the package declares no asset dependency.  Editorial formatting only, in addition: an AMS \texttt{amsart} preamble that restates the article's own declarations and macros used by the retained lines, namely the environments \texttt{proposition} on separate counters, together with the standard packages those lines need.  The retained environments print in this document as Proposition 1 in order of appearance, whereas the article prints the same items with the numbering of its own full document.  The complete native source of the article as distributed in the arXiv e-print for this exact version is bundled unchanged as \texttt{sources/208656/source0.tex}.
\begin{proposition}\label{prop:maximal-stratum-orbits}
Let \(1\le r\le n-1\). Then the set of \(G_0\)-orbits in
\(Y^0_{S,K_r}\) is naturally identified with $\binom{\mu_n}{r}\big/\mu_n,$
where $\mu_n=\{\zeta\in\mathbb C^*:\zeta^n=1\},$ and 
\(\binom{\mu_n}{r}\) denotes the set of \(r\)-element subsets of
\(\mu_n\), and \(\mu_n\) acts by multiplication.

Consequently,
\begin{equation}\label{eq:necklace-number}
N(n,r)
:=
\#\bigl(Y^0_{S,K_r}/G_0\bigr)
=
\frac1n
\sum_{d\mid\gcd(n,r)}
\varphi(d)
\binom{n/d}{r/d}.
\end{equation}
Moreover,
\begin{equation}\label{eq:single-orbit-normal}
N(n,r)=1
\quad\Longleftrightarrow\quad
r\in\{1,n-1\}
\quad\Longleftrightarrow\quad
K_r\in\mathcal N.
\end{equation}
\end{proposition}

\begin{proof}
Assume \(n\ge4\), and let \(1\le r\le n-1\). We put \(m=n-r\), and let $x=u(\mathbf a)w_0B/B.$
Suppose first that \(x\in Y^0_{S,K_r}\). Then
\(u(\mathbf a)w_0 \in B^-w_{K_r}B\). The first row of \(u(\mathbf a)w_0\) is
\[
(a_{n-1},a_{n-2},\ldots,a_1,1).
\]
Since the first block of \(w_{K_r}\) has size \(r\), the Bruhat
criterion gives that its first
nonzero entry must occur in column \(r\). As \(m=n-r\), this gives
\[
a_{m+1}=\cdots=a_{n-1}=0,
\qquad
a_m\neq0.
\]
We have,
\[
w_0\bigl(u(\mathbf a)w_0\bigr)^{-1}w_0
=
u(\mathbf a)^{-1}w_0.
\]
Moreover,
\[
u(\mathbf a)w_0 \in B^-w_{K_r}B
\quad\Longrightarrow\quad
w_0(u(\mathbf a)w_0)^{-1}w_0
\in
B^-\bigl(w_0w_{K_r}w_0\bigr)B
=
B^-w_{K_{n-r}}B.
\]
Writing $u(\mathbf a)^{-1}
=
I+\sum_{j=1}^{n-1}b_jN^j,$ the first row of \(u(\mathbf a)^{-1}w_0\) is $(b_{n-1},b_{n-2},\ldots,b_1,1).$
Since the first block of \(w_{K_{n-r}}\) has size \(n-r=m\), the
Bruhat criterion again  gives
\[
b_{r+1}=\cdots=b_{n-1}=0,
\qquad
b_r\neq0.
\]
Conversely, suppose these two sets of conditions hold, and let
\(J\subseteq S\) be the unique subset such that
\(x\in B^-w_JB/B\). The first set of conditions shows that the first
block associated with \(J\) has size \(r\), while the second shows
that its last block has size \(n-r\). Since these two sizes add up to
\(n\), the block decomposition is exactly \(n=r+(n-r)\). Hence
\(J=K_r\). Therefore
\begin{equation}\label{eq:maximal-stratum-toeplitz}
x\in Y^0_{S,K_r}
\quad\Longleftrightarrow\quad
\begin{cases}
a_{m+1}=\cdots=a_{n-1}=0,\quad a_m\neq0,\\
b_{r+1}=\cdots=b_{n-1}=0,\quad b_r\neq0.
\end{cases}
\end{equation}

Consider the polynomials
\[
A(T)=1+a_1T+\cdots+a_mT^m,
\qquad
B(T)=1+b_1T+\cdots+b_rT^r.
\]
Since $u(\mathbf a)u(\mathbf a)^{-1}=I$
and \(N^n=0\), we have 
\[
A(T)B(T)\equiv1\pmod{T^n}.
\]
Since $\deg A=m,\deg B=r$ and $ m+r=n$, 
we obtain
\begin{equation}\label{eq:AB-factorization}
A(T)B(T)=1+\gamma T^n
\end{equation}
for some $\gamma=a_mb_r\neq0.$

Suppose $B(T)$ factors as
\[
B(T)=\prod_{i=1}^r(1-z_iT).
\]
Since \(b_r\neq0\), all \(z_i\) are nonzero. By
\eqref{eq:AB-factorization}, we have $B(T)\mid 1+\gamma T^n.$
The polynomial \(1+\gamma T^n\) is separable, since
\(\gamma\neq0\). Hence \(B(T)\) is separable, so the \(z_i\)'s are
pairwise distinct. Substituting \(T=z_i^{-1}\) in
\eqref{eq:AB-factorization} gives
\[
z_i^n=-\gamma,
\qquad
1\le i\le r.
\]
Thus every point of \(Y^0_{S,K_r}\) determines an unordered set $\{z_1,\ldots,z_r\}
\subset\mathbb C^*$
of distinct elements having the same nonzero \(n\)-th power.

Conversely, suppose \(z_1,\ldots,z_r\) are distinct nonzero numbers
such that $z_1^n=\cdots=z_r^n=c\neq0.$
Then $B(T):=\prod_{i=1}^r(1-z_iT)$
divides \(1-cT^n\). Indeed, the roots of \(B(T)\) are the distinct numbers \(z_i^{-1}\),
and
\[
1-cz_i^{-n}=1-\frac{c}{z_i^n}=0.
\]
Hence every linear factor \(1-z_iT\) divides \(1-cT^n\), and since
these factors are pairwise coprime, their product \(B(T)\) divides
\(1-cT^n\). We put $A(T):=\frac{1-cT^n}{B(T)}.$
Then \(A(0)=1\), $\deg A=n-r=m,$ $\deg B=r$, and both leading coefficients are nonzero. Hence
\[
A(T)B(T)\equiv1\pmod{T^n},
\]
and \eqref{eq:maximal-stratum-toeplitz} gives a point of
\(Y^0_{S,K_r}\). We have therefore obtained a bijection between
\(Y^0_{S,K_r}\) and the unordered \(r\)-element subsets $\{z_1,\ldots,z_r\}\subset\mathbb C^*$
whose elements have a common nonzero \(n\)-th power.

We choose \(s\in\mathbb C^*\)
such that \(s^n=c\). Then
\[
C:=\{s^{-1}z_1,\ldots,s^{-1}z_r\}
\in\binom{\mu_n}{r},
\qquad
\{z_1,\ldots,z_r\}=sC.
\]
A different choice of \(s\) changes \(C\) by multiplication by an
element of \(\mu_n\). Under the \(G_0\)-action, the coefficients of
\[
B(T)=1+b_1T+\cdots+b_rT^r
\]
transform as \(b_j\mapsto t^{-j}b_j\). Thus \(B(T)\) transforms into
\[
1+\sum_{j=1}^r t^{-j}b_jT^j
=
B(t^{-1}T)
=
\prod_{i=1}^r\bigl(1-(t^{-1}z_i)T\bigr).
\]
Hence the associated subset transforms as $\{z_1,\ldots,z_r\}
\longmapsto
\{t^{-1}z_1,\ldots,t^{-1}z_r\}.$
Every \(G_0\)-orbit therefore has a representative in
\(\binom{\mu_n}{r}\), and two such representatives \(C,C'\) lie in
the same orbit if and only if $C'=\xi C$ for some $\xi\in\mu_n.$
Consequently,
\[
Y^0_{S,K_r}/G_0
\cong
\binom{\mu_n}{r}/\mu_n.
\]


% Choose \(s\in\mathbb C^*\) such that \(s^n=z_i^n\). Then $\{z_1,\ldots,z_r\}=sC$
% for some \(C\in\binom{\mu_n}{r}\). The \(G_0\)-action on the
% Toeplitz coordinates is
% \[
% t\cdot(a_1,\ldots,a_{n-1})
% =
% (t^{-1}a_1,\ldots,t^{-(n-1)}a_{n-1}),
% \]
% which corresponds to $\{z_1,\ldots,z_r\}
% \longmapsto
% \{t^{-1}z_1,\ldots,t^{-1}z_r\}.$
% Thus two subsets \(C,C'\subseteq\mu_n\) determine the same
% \(G_0\)-orbit precisely when $C'=\xi C$
% for some \(\xi\in\mu_n\). Hence
% \[
% Y^0_{S,K_r}/G_0
% \cong
% \binom{\mu_n}{r}/\mu_n.
% \]

The formula \eqref{eq:necklace-number} follows from Burnside's lemma.
Indeed, if \(\xi\in\mu_n\) has order \(d\), then its action on
\(\mu_n\) has \(n/d\) cycles, each of size \(d\). An \(r\)-element
subset is fixed by \(\xi\) if and only if \(d\mid r\), in which case
there are $\binom{n/d}{r/d}$
such subsets. Since there are \(\varphi(d)\) elements of order \(d\),
\eqref{eq:necklace-number} follows.

For \(r=1\), multiplication by \(\mu_n\) is transitive on
\(\mu_n\), and the case \(r=n-1\) follows by taking complements.
Hence \(N(n,r)=1\) for \(r=1,n-1\).

If \(2\le r\le n-2\), choose a primitive \(n\)-th root
\(\zeta\) and consider
\[
C_1=\{1,\zeta,\ldots,\zeta^{r-1}\},
\qquad
C_2=\{1,\zeta,\ldots,\zeta^{r-2},\zeta^r\}.
\]
The first contains \(r-1\) cyclically consecutive pairs, while the
second contains \(r-2\). This number is invariant under multiplication
by \(\mu_n\), so \(C_1\) and \(C_2\) belong to distinct
\(\mu_n\)-orbits. Thus \(N(n,r)\ge2\).

Finally, among the maximal subsets \(K_r\), we have
\[
K_r\in\mathcal N
\quad\Longleftrightarrow\quad
r\in\{1,n-1\},
\]
which proves \eqref{eq:single-orbit-normal}.
\end{proof}

\end{document}
